Par l'inégalité entre moyenne quadratique et moyenne arithmétique
(ou par convexité de $x \mapsto x^2$),
\[
\frac{a^2 + b^2 + c^2}{3} \geqslant \left( \frac{a+b+c}{3} \right)^2,
\qquad\text{d'où}\qquad
a^2 + b^2 + c^2 \geqslant \frac{(a+b+c)^2}{3}.
\]

Il suffit donc de montrer que $a + b + c \geqslant 3$. Or, par
l'inégalité arithmético-géométrique et l'hypothèse $abc = 1$ :
\[
\frac{a+b+c}{3} \geqslant \sqrt[3]{abc} = 1.
\]

En combinant les deux inégalités :
\[
a^2 + b^2 + c^2 \geqslant \frac{(a+b+c)^2}{3} = \frac{(a+b+c)(a+b+c)}{3}
\geqslant \frac{3\,(a+b+c)}{3} = a + b + c.
\]

L'égalité a lieu si et seulement si $a = b = c$ dans les deux
inégalités utilisées, c'est-à-dire $a = b = c = 1$ compte tenu de
$abc = 1$. \hfill $\blacksquare$
